\chapter{Rapidhash reference implementation}\label{sec:appendix_rapidhash}

This appendix gives the normative Python transcription of the rapidhash V3
topic hash function pinned by Cyphal v1.1
(section~\ref{sec:session_rapidhash}). It is bit-identical to the reference
C implementation on every platform and may be used directly to reproduce
the test vectors in section~\ref{sec:session_rapidhash}.
The C implementation is distributed under the MIT license by Nicolas De
Carli and is derived from the wyhash algorithm by Wang Yi.
All arithmetic is modulo $2^{64}$; intermediate results must be masked to
64 bits where the Python semantics do not do so automatically.

\begin{minted}[fontsize=\scriptsize]{python}
MASK64 = (1 << 64) - 1
SECRET = [ 0x2d358dccaa6c78a5, 0x8bb84b93962eacc9, 0x4b33a62ed433d4a3, 0x4d5a2da51de1aa47,
           0xa0761d6478bd642f, 0xe7037ed1a0b428db, 0x90ed1765281c388c, 0xaaaaaaaaaaaaaaaa]

def _read64(b, o): return int.from_bytes(b[o:o+8], "little")
def _read32(b, o): return int.from_bytes(b[o:o+4], "little")

def _mum(a, b):
    r = (a * b) & ((1 << 128) - 1)
    return (r & MASK64, (r >> 64) & MASK64)

def _mix(a, b):
    lo, hi = _mum(a, b)
    return lo ^ hi

def rapidhash(key: bytes, seed: int = 0) -> int:
    p, n = key, len(key)
    seed ^= _mix(seed ^ SECRET[2], SECRET[1])
    seed &= MASK64
    a = b = 0
    i = n
    if n <= 16:
        if n >= 4:
            seed ^= n
            if n >= 8:
                a = _read64(p, 0)
                b = _read64(p, n - 8)
            else:
                a = _read32(p, 0)
                b = _read32(p, n - 4)
        elif n > 0:
            a = (p[0] << 45) | p[n - 1]
            b = p[n >> 1]
    else:
        off = 0
        if i > 112:
            see1 = see2 = see3 = see4 = see5 = see6 = seed
            while True:
                seed = _mix(_read64(p, off)      ^ SECRET[0], _read64(p, off +  8) ^ seed)
                see1 = _mix(_read64(p, off + 16) ^ SECRET[1], _read64(p, off + 24) ^ see1)
                see2 = _mix(_read64(p, off + 32) ^ SECRET[2], _read64(p, off + 40) ^ see2)
                see3 = _mix(_read64(p, off + 48) ^ SECRET[3], _read64(p, off + 56) ^ see3)
                see4 = _mix(_read64(p, off + 64) ^ SECRET[4], _read64(p, off + 72) ^ see4)
                see5 = _mix(_read64(p, off + 80) ^ SECRET[5], _read64(p, off + 88) ^ see5)
                see6 = _mix(_read64(p, off + 96) ^ SECRET[6], _read64(p, off +104) ^ see6)
                off += 112
                i   -= 112
                if i <= 112:
                    break
            seed ^= see1
            see2 ^= see3
            see4 ^= see5
            seed ^= see6
            see2 ^= see4
            seed ^= see2
            seed &= MASK64
        if i > 16:
            seed = _mix(_read64(p, off) ^ SECRET[2], _read64(p, off +  8) ^ seed)
            if i > 32:
                seed = _mix(_read64(p, off + 16) ^ SECRET[2], _read64(p, off + 24) ^ seed)
                if i > 48:
                    seed = _mix(_read64(p, off + 32) ^ SECRET[1], _read64(p, off + 40) ^ seed)
                    if i > 64:
                        seed = _mix(_read64(p, off + 48) ^ SECRET[1], _read64(p, off + 56) ^ seed)
                        if i > 80:
                            seed = _mix(_read64(p, off + 64) ^ SECRET[2], _read64(p, off + 72) ^ seed)
                            if i > 96:
                                seed = _mix(_read64(p, off + 80) ^ SECRET[1], _read64(p, off + 88) ^ seed)
        a = _read64(p, off + i - 16) ^ i
        b = _read64(p, off + i - 8)
    a ^= SECRET[1]
    b ^= seed
    lo, hi = _mum(a, b)
    return _mix(lo ^ SECRET[7], hi ^ SECRET[1] ^ i)
\end{minted}
